\chapter{Modèles continus et méthode d'Euler}
\section*{Exercice 13.1 -- Croissance continue}
Résoudre $y'=0{,}2y$, $y(0)=50$. Quel est le taux d'évolution exact sur une unité de temps ?

\section*{Exercice 13.2 -- Décroissance et demi-vie}
Pour $y(t)=80e^{-0{,}4t}$, exprimer le temps où $y(t)=40$. Le logarithme peut être utilisé comme complément.

\section*{Exercice 13.3 -- Apport constant}
Résoudre $y'=-2y+6$, $y(0)=1$, en déplaçant l'équilibre.

\section*{Exercice 13.4 -- Deux pas d'Euler}
Approcher $y(1)$ pour $y'=y$, $y(0)=1$, avec deux pas égaux. Comparer à $e$.

\section*{Exercice 13.5 -- Un pas inadéquat}
Euler appliqué à $y'=-3y$, $y(0)=1$, avec $h=1$ donne quelle suite ? Est-elle compatible avec la solution exacte ?

\section*{Exercice 13.6 -- Unités du modèle}
Dans $y'=ay+b$, $y$ est mesuré en litres et $t$ en minutes. Quelles sont les unités de $a,b$ ? Que devient $a$ si l'on mesure le temps en secondes ?
